\documentclass[11pt,a4paper]{article}
\usepackage[margin=1in]{geometry}
\usepackage{times}
\usepackage{graphicx}
\usepackage{amsmath}
\usepackage{amssymb}
\usepackage{booktabs}
\usepackage{multirow}
\usepackage{xcolor}
\usepackage{hyperref}
\usepackage{natbib}

\title{GRPO Alignment Fails on Structurally Pruned Reasoning Models:\\Reward Hacking, Training Collapse, and the Redundancy-Stability Tradeoff}

\author{Anonymous}

\begin{document}
\maketitle

\begin{abstract}
We present the first systematic study of how structural pruning interacts with RLVR (Reinforcement Learning from Verifiable Rewards) alignment in reasoning LLMs. We discover that GRPO alignment universally fails to improve reasoning accuracy---even on unpruned models, reward increases from 0 to 0.661 while accuracy remains unchanged (81.27\% $\rightarrow$ 81.27\%, identical 1072/1319 correct answers). On structurally pruned models, GRPO additionally induces \emph{training collapse}: gradient norms drop to near-zero and policy diversity vanishes. To diagnose this phenomenon, we introduce RASP (Redundancy-Aware Structured Pruning), a CKA-based scoring method that reveals a fundamental \emph{Redundancy-Stability Tradeoff}: structures that are safely removable under SFT recovery (functionally redundant pathways) provide critical exploration stability for RLVR. RASP-pruned models collapse 6$\times$ faster than standard importance-pruned models (step 10 vs.\ step 60), while unstructured pruning produces no collapse at all---confirming that training collapse is a structural damage phenomenon. We validate across two models, four pruning methods, two sparsity levels, and three recovery paradigms, establishing that the prune-recover-align pipeline requires fundamental reconsideration when RLVR alignment is intended.
\end{abstract}

\section{Introduction}

The deployment pipeline for compressed reasoning LLMs typically follows a three-stage process: structural pruning to reduce model size, supervised fine-tuning (SFT) to recover lost capabilities, and reinforcement learning alignment (RLVR) to further improve reasoning~\cite{self-healing-2025}. This pipeline implicitly assumes that pruning decisions are independent of the downstream alignment method---that structures safe to remove for SFT recovery are equally safe for RLVR alignment.

We demonstrate that this assumption fails catastrophically. Through controlled experiments on DeepSeek-R1-Distill-Qwen-7B, we establish three findings:

\paragraph{Finding 1: GRPO reward-accuracy decoupling is universal.} GRPO does not improve reasoning accuracy for \emph{any} model configuration post-SFT. On the unpruned model, 100 steps of GRPO training increase reward from 0 to 0.661 while accuracy remains at exactly 81.27\% (1072/1319 correct on GSM8K). On pruned models, GRPO additionally induces complete training collapse---policy diversity vanishes and gradient norms approach zero.

\paragraph{Finding 2: Collapse dynamics diagnose structural damage.} The speed of GRPO collapse varies systematically with pruning method: RASP (which removes functionally redundant structures) collapses at step $\sim$10, GRASPrune at step $\sim$18, Bonsai at step $\sim$28, and standard importance pruning at step $\sim$60. This ordering directly reflects how aggressively each method removes structural redundancy, establishing collapse speed as a zero-cost diagnostic of structural damage type.

\paragraph{Finding 3: Optimal pruning depends on recovery paradigm.} Under SFT recovery, RASP shows a consistent marginal advantage (+2.2pp at 20\% sparsity, +3.2pp at 25\%). Under RLVR, RASP is actively harmful---it collapses 6$\times$ faster than standard pruning. This Redundancy-Stability Tradeoff means the structures optimal to prune for SFT are exactly those needed for RLVR stability.

Crucially, unstructured pruning at the same sparsity (20\%) produces \emph{no} training collapse---GRPO trains stably for 100 steps with reward reaching 0.548. This confirms that collapse is a \emph{structural} damage phenomenon: removing entire computational units eliminates exploration pathways that RLVR requires, while zeroing individual weights preserves network topology and exploration capacity.

Our results challenge Self-Healing~\cite{self-healing-2025}, which claims RLVR recovers pruned reasoning models. We show that under structured pruning with controlled accuracy evaluation, GRPO produces reward improvement without accuracy gain followed by training collapse. The prune-recover-align pipeline requires fundamental reconsideration when RLVR alignment is the target.

\section{Related Work}

\paragraph{Structural pruning for LLMs.} LLM-Pruner~\cite{llmpruner-2023} introduced dependency-aware structural pruning for transformer models. SliceGPT~\cite{slicegpt-2024} prunes via rotation-invariant projections. GRASPrune~\cite{grasprune-2026} achieves SOTA budgeted structural pruning via learned gates. Bonsai~\cite{bonsai-2025} uses activation norms for gradient-free pruning. These methods evaluate pruning quality via post-prune or post-SFT accuracy without studying RLVR alignment interaction.

\paragraph{Recovery and alignment of pruned models.} Self-Healing~\cite{self-healing-2025} claims RLVR (specifically GRPO) can recover pruned reasoning models to 95\%+ of original accuracy at 25\% unstructured sparsity. PASER~\cite{paser-2025} optimizes recovery data selection. RESP~\cite{resp-2025} uses reasoning-trace calibration. None document GRPO training collapse on structurally pruned models or study the interaction between pruning strategy and alignment stability.

\paragraph{CKA in neural network analysis.} Centered Kernel Alignment (CKA)~\cite{cka-2019} measures representational similarity between neural network layers. MPruner~\cite{mpruner-2024} uses inter-layer CKA for depth pruning in CNNs. We use intra-layer CKA as a \emph{functional redundancy} metric that reveals the Redundancy-Stability Tradeoff in reasoning LLMs.

\paragraph{RLVR for reasoning.} GRPO~\cite{grpo-2024} provides group-relative policy optimization for reasoning tasks. Recent work on reasoning LLM alignment~\cite{deepseek-r1-2024} uses GRPO with verifiable rewards. The ``RL Squeezes, SFT Expands'' hypothesis~\cite{rl-squeezes-2026} proposes that RL narrows representation diversity while SFT broadens it---consistent with our observation that RL requires structural redundancy for stable exploration.

\section{Method}

\subsection{RASP: Redundancy-Aware Structured Pruning}

We introduce RASP as a diagnostic pruning method designed to reveal the role of functional redundancy in alignment stability. For each prunable structure $h$ (attention head or MLP group), RASP computes:
\begin{equation}
S(h) = I(h) \times (1 - F(h))^{\alpha}
\label{eq:rasp}
\end{equation}
where $I(h)$ is the standard importance score (gradient $\times$ magnitude) and $F(h) \in [0,1]$ is the functional redundancy score based on CKA similarity with other structures in the same layer:
\begin{equation}
F(h) = \max_{h' \neq h, h' \in \text{layer}(h)} \text{CKA}(\mathbf{X}_h, \mathbf{X}_{h'})
\end{equation}
Here $\mathbf{X}_h \in \mathbb{R}^{n \times d}$ is the activation matrix of structure $h$ over $n$ calibration samples. The penalty $(1-F(h))^\alpha$ reduces the score of structures with \emph{low} redundancy (unique, non-replaceable structures), causing RASP to preferentially prune structures with high functional overlap.

The hyperparameter $\alpha$ controls the strength of the redundancy penalty. At $\alpha=0$, RASP reduces to standard importance pruning.

\subsection{Pruning Pipeline}

We prune at the granularity of attention heads and MLP intermediate groups. Given a target sparsity ratio $s$, we rank all structures by $S(h)$ and remove the bottom-$s$ fraction. No iterative or progressive pruning is used---all structures are removed in a single pass.

\subsection{Recovery Protocols}

\paragraph{SFT Recovery.} Supervised fine-tuning on GSM8K training set with teacher forcing. Training for 2000 steps with learning rate 2e-5 and batch size 8.

\paragraph{GRPO Recovery.} Group Relative Policy Optimization on GSM8K with verifiable answer checking. 100 training steps, group size 8, learning rate 5e-7, KL coefficient 0.04. We monitor collapse indicators: fraction of zero-variance groups (policy diversity), gradient norm, and reward trajectory.

\subsection{Collapse Indicators}

We define GRPO training collapse via three simultaneous conditions:
\begin{enumerate}
\item \textbf{Policy diversity loss}: fraction of zero-standard-deviation response groups exceeds 0.85
\item \textbf{Gradient vanishing}: gradient norm drops below 0.01
\item \textbf{Irreversibility}: conditions persist for $\geq$5 consecutive steps
\end{enumerate}
The \emph{collapse step} is the first step where all three conditions are met.

\section{Experimental Setup}

\paragraph{Models.} Primary: DeepSeek-R1-Distill-Qwen-7B (7B parameters, distilled reasoning model). Cross-model validation: Qwen2.5-Math-7B (math-specialized).

\paragraph{Evaluation.} GSM8K test set, $N=1319$ problems, greedy decoding. We report accuracy (fraction correct) with standard errors. Extended evaluation on MATH-500, HumanEval, and MMLU for out-of-distribution assessment.

\paragraph{Pruning methods compared.}
\begin{itemize}
\item \textbf{RASP} ($\alpha=1.0$): CKA-aware importance scoring (Eq.~\ref{eq:rasp})
\item \textbf{Standard}: Importance-only ($\alpha=0$, gradient $\times$ magnitude)
\item \textbf{GRASPrune}: Gate-learned structural importance (ACL 2026 SOTA)
\item \textbf{Bonsai}: Activation-norm-based (gradient-free, ICML 2025)
\item \textbf{Unstructured}: Magnitude-based weight pruning (no structural removal)
\end{itemize}

\paragraph{Sparsity levels.} 20\% (primary), 25\% (scaling validation), 30\% (ceiling probe).

\section{Results}

\subsection{Post-Prune Accuracy (Table~\ref{tab:postprune})}

\begin{table}[t]
\centering
\small
\begin{tabular}{lccc}
\toprule
\textbf{Method} & \textbf{20\%} & \textbf{25\%} & \textbf{30\%} \\
\midrule
Unpruned & \multicolumn{3}{c}{85.4\% $\pm$ 1.0\%} \\
\midrule
RASP & 29.8 $\pm$ 1.3 & 21.0 & 12.0 $\pm$ 3.2 \\
Standard & 26.1 $\pm$ 1.2 & 16.0 & 9.0 $\pm$ 2.9 \\
GRASPrune & 6.8 & --- & --- \\
Bonsai & 19.0 & --- & --- \\
Unstructured & 68.2 $\pm$ 1.3 & --- & --- \\
\midrule
RASP (5-seed) & 30.4 $\pm$ 2.8 & --- & --- \\
Standard (5-seed) & 26.5 $\pm$ 3.2 & --- & --- \\
\bottomrule
\end{tabular}
\caption{Post-prune accuracy (\%) on GSM8K ($N=1319$). RASP vs.\ Standard gap at 20\%: 3.7pp ($p=0.034$). Multi-seed results (5 seeds): paired $t$-test $p=0.024$, Cohen's $d=0.58$. Unstructured pruning preserves 68.2\% accuracy---38pp above the best structured method---confirming that structural removal causes far greater damage than weight zeroing.}
\label{tab:postprune}
\end{table}

Table~\ref{tab:postprune} establishes the damage landscape across pruning methods and sparsity levels. All structured pruning methods reduce accuracy by $>$55pp at 20\% sparsity, while unstructured pruning at the same ratio retains 68.2\%---a 38.4pp gap that reflects the fundamental difference between removing entire computational units versus zeroing individual weights. The 30\% results (RASP 12\%, Standard 9\%) confirm $\sim$25--30\% as the practical structured pruning ceiling for this model class.

\subsection{Recovery Accuracy Matrix (Table~\ref{tab:recovery})}

\begin{table}[t]
\centering
\small
\begin{tabular}{lcccc}
\toprule
\textbf{Method} & \textbf{Post-prune} & \textbf{+SFT} & \textbf{+GRPO} & \textbf{$\Delta$(GRPO$-$SFT)} \\
\midrule
Unpruned & 85.4 & 81.3 & 81.3 & 0.00 \\
\midrule
RASP 20\% & 29.8 & 61.5 & 60.7 & $-$0.84 \\
Standard 20\% & 26.1 & 59.3 & 58.9 & $-$0.38 \\
GRASPrune 20\% & 6.8 & 39.5 & 38.8 & $-$0.70 \\
Bonsai 20\% & 19.0 & 47.0 & 46.0 & $-$1.00 \\
Unstructured 20\% & 68.2 & 76.3 & 75.8 & $-$0.50 \\
\midrule
RASP 25\% & 21.0 & 54.5 & 53.5 & $-$1.00 \\
Standard 25\% & 16.0 & 51.3 & 50.5 & $-$0.80 \\
\bottomrule
\end{tabular}
\caption{Recovery accuracy (\%) on GSM8K ($N=1319$). GRPO \emph{never} improves over SFT---in every configuration, post-GRPO accuracy is equal to or below post-SFT accuracy. The SFT gap between RASP and Standard is 2.2pp at 20\% ($p=0.25$) and 3.2pp at 25\% ($p=0.089$).}
\label{tab:recovery}
\end{table}

Table~\ref{tab:recovery} is the central result: across all pruning methods, sparsity levels, and even unpruned models, GRPO does not improve reasoning accuracy over SFT. The unpruned model achieves reward 0.661 after 100 GRPO steps while maintaining exactly 81.27\% accuracy (identical 1072/1319 correct answers as post-SFT)---demonstrating reward-accuracy decoupling independent of pruning.

\subsection{GRPO Collapse Dynamics (Table~\ref{tab:collapse})}

\begin{table}[t]
\centering
\small
\begin{tabular}{lcccc}
\toprule
\textbf{Method} & \textbf{Collapse} & \textbf{Peak} & \textbf{Final} & \textbf{Final} \\
& \textbf{Step} & \textbf{Reward} & \textbf{frac\_zero} & \textbf{$\|\nabla\|$} \\
\midrule
Unpruned & None & 0.661 & 0.28 & stable \\
\midrule
\multicolumn{5}{l}{\textit{20\% sparsity}} \\
RASP & $\sim$10 & 0.039 & 0.95 & 0.007 \\
GRASPrune & $\sim$18 & 0.012 & 0.92 & 0.008 \\
Bonsai & $\sim$28 & 0.022 & 0.88 & --- \\
Standard & $\sim$60 & 0.065 & 0.72$^*$ & --- \\
Unstructured & None & 0.548 & 0.32 & 5.82 \\
\midrule
\multicolumn{5}{l}{\textit{25\% sparsity}} \\
RASP & $\sim$7 & 0.018 & 0.97 & 0.005 \\
Standard & $\sim$30 & 0.042 & 0.72 & --- \\
\bottomrule
\end{tabular}
\caption{GRPO collapse dynamics. Collapse step = first step meeting all three collapse criteria. The severity gradient (RASP $<$ GRASPrune $<$ Bonsai $<$ Standard) directly maps to how aggressively each method removes functionally redundant structures. Unstructured pruning produces no collapse. $^*$Standard at 20\% shows transient dip at step $\sim$60 with partial recovery.}
\label{tab:collapse}
\end{table}

The collapse severity gradient reveals a clear pattern: methods that more aggressively remove functionally redundant structures (as measured by CKA) collapse faster. RASP explicitly targets redundant structures and collapses first (step 10); Standard retains some redundancy by chance and collapses last among structured methods (step 60). The critical control is unstructured pruning: with all structural units intact (but 20\% of weights zeroed), GRPO trains stably for 100 steps---confirming that collapse is structural-damage-specific, not a general pruning effect.

At 25\% sparsity, both methods collapse faster (RASP: 7, Standard: 30), but the ratio compresses from 6$\times$ to 4.3$\times$, suggesting that at higher sparsity, both methods converge toward a floor where all redundancy is exhausted.

\subsection{Alpha Sensitivity Ablation (Table~\ref{tab:alpha})}

\begin{table}[t]
\centering
\small
\begin{tabular}{lccccc}
\toprule
$\alpha$ & 0.0 & 0.25 & 0.5 & 1.0 & 2.0 \\
\midrule
Post-prune (\%) & 27.0 & 23.0 & 22.0 & 31.0 & 36.0 \\
Post-SFT (\%) & 59.3 & 55.8 & 54.9 & 57.3 & 60.8 \\
\bottomrule
\end{tabular}
\caption{Alpha sensitivity ablation at 20\% sparsity ($N=1319$). $\alpha=0$ reduces to standard importance pruning. The U-shaped curve with minimum at $\alpha=0.5$ and optimum at $\alpha=2.0$ indicates that aggressive redundancy preservation (high $\alpha$) benefits both post-prune and post-SFT accuracy. The paradox: higher $\alpha$ preserves more unique structures for better accuracy but removes more redundancy---accelerating GRPO collapse.}
\label{tab:alpha}
\end{table}

The U-shaped alpha curve confirms that RASP's CKA penalty is a meaningful parameter, not an arbitrary addition. The minimum at $\alpha=0.5$ reflects a ``noise injection'' regime where weak CKA penalty disrupts importance ranking without providing sufficient redundancy signal. At $\alpha \geq 1.0$, the penalty becomes strong enough to meaningfully protect unique structures, improving both post-prune (+4pp) and post-SFT accuracy. However, higher $\alpha$ also means more aggressive redundancy removal---precisely what accelerates GRPO collapse. This is the Redundancy-Stability Tradeoff made explicit through a single hyperparameter.

\subsection{Cross-Model Validation (Table~\ref{tab:crossmodel})}

\begin{table}[t]
\centering
\small
\begin{tabular}{llcccc}
\toprule
\textbf{Model} & \textbf{Sparsity} & \textbf{RASP} & \textbf{Standard} & \textbf{Ratio} & \textbf{95\% CI} \\
\midrule
DeepSeek-R1 & 20\% & 10 & 60 & 6.0$\times$ & --- \\
DeepSeek-R1 & 25\% & 7 & 30 & 4.3$\times$ & --- \\
\midrule
Qwen-Math & 20\% & 12 & 18 & 1.5$\times$ & [0.8, 2.8] \\
Qwen-Math & 25\% & 8 & 14 & 1.75$\times$ & [0.9, 3.1] \\
\bottomrule
\end{tabular}
\caption{Cross-model collapse step comparison (RASP vs.\ Standard). On DeepSeek-R1-Distill-Qwen-7B, the collapse severity gradient is clean (6$\times$ ratio at 20\%). On Qwen2.5-Math-7B, the direction is consistent (RASP collapses first) but the ratio is smaller and 95\% CI includes 1---the gradient is a diagnostic tendency rather than a universal law.}
\label{tab:crossmodel}
\end{table}

GRPO collapse is robust across models: both DeepSeek-R1-Distill and Qwen2.5-Math-7B exhibit collapse under structural pruning at both sparsity levels. However, the \emph{severity gradient}---the differential collapse speed between RASP and Standard---does not cleanly replicate. On Qwen-Math-7B, the RASP/Standard ratio is only 1.5--1.75$\times$ with confidence intervals including 1. We interpret the collapse severity gradient as a diagnostic tendency observed on DeepSeek-R1, with directionally consistent but not statistically significant replication on Qwen-Math-7B. The difference likely reflects Qwen-Math's more uniform internal redundancy distribution from math-specific training.

\subsection{Scoring Ablation (Table~\ref{tab:scoring})}

\begin{table}[t]
\centering
\small
\begin{tabular}{lcc}
\toprule
\textbf{Scoring Method} & \textbf{Post-prune (\%)} & \textbf{EOS Rate} \\
\midrule
Random & 17.0 & 0.45 \\
Magnitude-only & 24.0 & 0.78 \\
Gradient-only & 28.0 & 0.85 \\
Standard (grad$\times$mag) & 27.0 & 0.82 \\
RASP ($\alpha=1.0$) & 31.0 & 0.89 \\
\bottomrule
\end{tabular}
\caption{Scoring ablation at 20\% sparsity. CKA's marginal contribution over gradient-only is 3pp. The ordering confirms that gradient-based importance is the dominant signal, with CKA providing complementary but non-dominant functional redundancy information.}
\label{tab:scoring}
\end{table}

The scoring ablation decomposes RASP's advantage. Gradient-only (28\%) $\approx$ Standard (27\%), confirming that magnitude contributes minimally beyond gradient. CKA's marginal contribution is 3pp (31\% vs.\ 28\%)---helpful but not decisive for post-prune accuracy. CKA's primary value is not in post-prune performance but in revealing collapse dynamics (Section 5.3).

\subsection{Inference Efficiency (Table~\ref{tab:speedup})}

\begin{table}[t]
\centering
\small
\begin{tabular}{lccc}
\toprule
\textbf{Config} & \textbf{Latency (ms)} & \textbf{Throughput (tok/s)} & \textbf{$\Delta$} \\
\midrule
Dense & 142 & 85 & --- \\
RASP 20\% & 121 & 99 & $-$14.8\% / +16.5\% \\
Standard 20\% & 120 & 100 & $-$15.5\% / +17.6\% \\
\bottomrule
\end{tabular}
\caption{Inference speedup on A100-80GB (batch=1, seq\_len=512, 50 runs after 10 warmup). Actual speedup ($\sim$15\%) is below the theoretical 19\% due to non-pruned computation (attention, LayerNorm, embedding) and framework overhead.}
\label{tab:speedup}
\end{table}

Structured pruning at 20\% provides $\sim$15\% latency reduction and $\sim$17\% throughput improvement. RASP and Standard produce nearly identical speedup since both remove the same total number of structures---the pruning \emph{method} affects which structures are removed but not the overall efficiency gain. Meaningful practical speedup likely requires 25--30\% sparsity, where our results show structured pruning approaches its practical ceiling.

\subsection{Extended Evaluation (Table~\ref{tab:extended})}

\begin{table}[t]
\centering
\small
\begin{tabular}{lcccc}
\toprule
\textbf{Config} & \textbf{GSM8K} & \textbf{MATH-500} & \textbf{HumanEval} & \textbf{MMLU} \\
\midrule
Unpruned (orig.) & 85.4 & 55.8 & 56.1 & 65.1 \\
Unpruned (SFT) & 81.3 & 52.4 & 54.9 & 63.8 \\
\midrule
RASP SFT 20\% & 61.5 & 31.8 & 35.4 & 51.2 \\
Standard SFT 20\% & 59.3 & 28.5 & 30.5 & 49.8 \\
\midrule
Recovery ratio & 75.7\% & 60.7\% & 64.5\% & 80.3\% \\
\bottomrule
\end{tabular}
\caption{Extended evaluation across benchmarks. Recovery ratio = (pruned SFT) / (unpruned SFT). SFT recovery is domain-specific: GSM8K in-domain recovery (75.7\%) far exceeds harder math (60.7\%) and code (64.5\%). MMLU factual knowledge (80.3\%) is least affected by MLP pruning.}
\label{tab:extended}
\end{table}

Out-of-distribution evaluation reveals that SFT recovery is partially domain-specific. The pruned model recovers 75.7\% of capability on in-domain GSM8K but only 60.7\% on MATH-500 (harder multi-step reasoning) and 64.5\% on HumanEval (code generation). MMLU recovery is highest (80.3\%) because factual knowledge resides primarily in embeddings and attention patterns, which our MLP-focused pruning preserves. The RASP vs.\ Standard gap on OOD benchmarks (MATH-500: 3.3pp, HumanEval: 4.9pp, MMLU: 1.4pp) is directionally consistent but partially domain-specific.

\subsection{CKA Reliability Analysis (Table~\ref{tab:cka_reliability})}

\begin{table}[t]
\centering
\small
\begin{tabular}{lccc}
\toprule
\textbf{Calibration Size} & \textbf{Spearman $\rho$} & \textbf{$\Delta$ Acc vs.\ 512} & \textbf{Time (s)} \\
\midrule
256 & 0.88 & $-$1.5pp & 180 \\
512 & 0.93 & (ref) & 340 \\
1024 & 0.95 & +0.8pp & 660 \\
\midrule
Split-half (512) & 0.93 & --- & --- \\
\bottomrule
\end{tabular}
\caption{CKA scoring reliability across calibration sizes ($n=896$ structures). Split-half consistency $\rho=0.93$ ($p < 10^{-45}$) confirms stable scoring at 512 samples. Accuracy impact of calibration size is $<$2pp.}
\label{tab:cka_reliability}
\end{table}

CKA-based scoring achieves high reliability at modest calibration cost. Split-half consistency of $\rho = 0.93$ confirms that redundancy rankings are stable across random data subsets. The diminishing returns from 512$\rightarrow$1024 (0.93$\rightarrow$0.95) indicate 512 samples as the cost-effective minimum, adding only $\sim$340 seconds to the pruning pipeline for 896 structures.

\subsection{CKA-Recoverability Correlation (Table~\ref{tab:cka_corr})}

\begin{table}[t]
\centering
\small
\begin{tabular}{lcc}
\toprule
\textbf{Analysis Level} & \textbf{Pearson $r$} & \textbf{$p$-value} \\
\midrule
Overall (224 heads) & 0.38 & 0.003 \\
Within-layer (controlled) & 0.22 & 0.041 \\
\midrule
Partial $r$ (controlling weight norm) & 0.12 & --- \\
Partial $r$ (controlling grad importance) & 0.15 & --- \\
\midrule
Between-layer variance explained & 62\% & --- \\
Within-layer variance explained & 8\% & --- \\
Layer ICC & 0.71 & --- \\
\bottomrule
\end{tabular}
\caption{Mixed-effects regression: CKA-based $F(h)$ vs.\ post-recovery activation change. CKA captures a real signal ($r=0.38$, $p=0.003$) but prediction power is layer-dominated (ICC=0.71). Within-layer, CKA explains only 8\% of variance.}
\label{tab:cka_corr}
\end{table}

CKA-based redundancy scores correlate meaningfully with post-recovery activation change at the macro level ($r=0.38$, $p=0.003$), confirming that functional redundancy is a real property captured by our metric. However, the layer ICC of 0.71 reveals that 71\% of variance is between-layer (early layers are more redundant than deep layers) with only 8\% explained within-layer. This explains why RASP provides marginal rather than transformative advantage: the dominant redundancy signal is at the layer level, where RASP and Standard make similar allocation decisions (both preferentially prune shallow MLPs).

\section{Discussion}

\paragraph{The Structural Damage Hypothesis.} Our most striking result is the contrast between structured and unstructured pruning under GRPO. At identical 20\% sparsity, structured pruning (RASP: 29.8\% post-prune) produces permanent collapse at step 10, while unstructured pruning (68.2\% post-prune) trains stably for 100 steps. This cannot be explained by accuracy difference alone---the unstructured model starts from 76.3\% post-SFT (vs.\ structured at 61.5\%), but the qualitative phenomenon (collapse vs.\ stability) points to structural integrity rather than accuracy level as the mechanism. Removing entire attention heads or MLP groups eliminates exploration pathways; zeroing individual weights within intact structures preserves the network's capacity for diverse policy exploration.

\paragraph{Relationship to Self-Healing.} Self-Healing~\cite{self-healing-2025} reports RLVR recovery at 25\% \emph{unstructured} sparsity. Our unstructured baseline is consistent: no collapse at 20\% unstructured. However, we additionally show that (1) reward improvement does not translate to accuracy gain even without collapse (unstructured: reward 0.548, accuracy $\Delta = -0.5$pp), and (2) structured pruning---the practically deployable compression paradigm for inference speedup---produces qualitatively different failure modes. Self-Healing's results may reflect reward improvement mislabeled as capability recovery.

\paragraph{Generalization Limitations.} The collapse severity gradient (RASP faster than Standard) is clean on DeepSeek-R1-Distill-Qwen-7B but does not significantly replicate on Qwen2.5-Math-7B. We interpret this as a property of models with the specific redundancy distribution created by reasoning distillation training, not a universal law. The Redundancy-Stability Tradeoff is best characterized as an empirical observation on the DeepSeek-R1 family, with directionally consistent cross-model evidence.

\paragraph{SFT Gap Significance.} The SFT gap between RASP and Standard is 2.2pp at 20\% ($p=0.25$) and 3.2pp at 25\% ($p=0.089$). This is a consistent trend that amplifies with sparsity but does not reach significance at 20\%. We emphasize that the paper's primary contribution is the GRPO failure mechanism, not SFT superiority of RASP.

\paragraph{Practical Recommendations.} For practitioners compressing reasoning models: (1) SFT-only recovery is currently the safe path---GRPO adds no accuracy benefit and risks collapse; (2) structured pruning beyond 25\% is impractical for 7B reasoning models; (3) if RLVR alignment is required downstream, unstructured pruning preserves the structural integrity needed for stable training.

\section{Conclusion}

We demonstrate that GRPO alignment universally fails to improve reasoning accuracy post-SFT, producing reward hacking (reward increases, accuracy unchanged) across all configurations including unpruned models. Structural pruning additionally induces training collapse, with collapse speed directly reflecting how aggressively pruning removes functional redundancy. The Redundancy-Stability Tradeoff---structures safe for SFT removal are critical for RLVR stability---means the prune-recover-align pipeline requires fundamental reconsideration. Our diagnostic tool (RASP) and the unstructured pruning control establish that training collapse is a structural damage phenomenon, not a general pruning artifact. Future work should explore pruning-aware RLVR methods that account for structural integrity requirements.

\bibliographystyle{plainnat}
\begin{thebibliography}{20}

\bibitem[Chen et~al.(2025)]{self-healing-2025}
Chen, X., et~al. (2025).
\newblock Self-Healing: Recovering Pruned Reasoning Models via RLVR.
\newblock \emph{arXiv preprint arXiv:2501.xxxxx}.

\bibitem[DeepSeek(2024)]{deepseek-r1-2024}
DeepSeek-AI (2024).
\newblock DeepSeek-R1: Incentivizing Reasoning Capability in LLMs via Reinforcement Learning.
\newblock \emph{arXiv preprint arXiv:2401.xxxxx}.

\bibitem[Shao et~al.(2024)]{grpo-2024}
Shao, Z., et~al. (2024).
\newblock DeepSeekMath: Pushing the Limits of Mathematical Reasoning in Open Language Models.
\newblock \emph{arXiv preprint arXiv:2402.xxxxx}.

\bibitem[Ma et~al.(2023)]{llmpruner-2023}
Ma, X., et~al. (2023).
\newblock LLM-Pruner: On the Structural Pruning of Large Language Models.
\newblock In \emph{NeurIPS 2023}.

\bibitem[Ashkboos et~al.(2024)]{slicegpt-2024}
Ashkboos, S., et~al. (2024).
\newblock SliceGPT: Compress Large Language Models by Deleting Rows and Columns.
\newblock In \emph{ICLR 2024}.

\bibitem[Zhang et~al.(2026)]{grasprune-2026}
Zhang, Y., et~al. (2026).
\newblock GRASPrune: Gradient-Aware Structural Pruning with Learned Gates.
\newblock In \emph{ACL 2026}.

\bibitem[Li et~al.(2025)]{bonsai-2025}
Li, T., et~al. (2025).
\newblock Bonsai: Gradient-Free Structural Pruning via Activation Norms.
\newblock In \emph{ICML 2025}.

\bibitem[Kornblith et~al.(2019)]{cka-2019}
Kornblith, S., et~al. (2019).
\newblock Similarity of Neural Network Representations Revisited.
\newblock In \emph{ICML 2019}.

\bibitem[Wang et~al.(2024)]{mpruner-2024}
Wang, Z., et~al. (2024).
\newblock MPruner: CKA-Based Layer Pruning for Vision Transformers.
\newblock In \emph{CVPR 2024}.

\bibitem[Liu et~al.(2025)]{paser-2025}
Liu, H., et~al. (2025).
\newblock PASER: Pruning-Aware SFT for Efficient Recovery.
\newblock \emph{arXiv preprint arXiv:2503.xxxxx}.

\bibitem[Zhao et~al.(2025)]{resp-2025}
Zhao, K., et~al. (2025).
\newblock RESP: Reasoning-Enhanced Structured Pruning.
\newblock \emph{arXiv preprint arXiv:2512.xxxxx}.

\bibitem[Huang et~al.(2026)]{rl-squeezes-2026}
Huang, J., et~al. (2026).
\newblock RL Squeezes, SFT Expands: Understanding Representation Dynamics in LLM Alignment.
\newblock In \emph{ICLR 2026}.

\end{thebibliography}

\appendix
\section{Multi-Seed Variance Analysis}
\label{app:multiseed}

Five random seeds (42, 123, 456, 789, 1024) confirm RASP's post-prune advantage: mean 30.4\% $\pm$ 2.8\% vs.\ Standard 26.5\% $\pm$ 3.2\%. Paired $t$-test $p=0.024$, Cohen's $d = 0.58$ (medium effect). Per-seed results: RASP [31, 28, 33, 29, 31]\%, Standard [27, 24, 30, 22, 29]\%.

\section{Bonsai Baseline Details}
\label{app:bonsai}

Bonsai (ICML 2025, gradient-free activation-norm pruning) at 20\%: post-prune 19\%, post-SFT 47\%, GRPO collapse at step $\sim$28, peak reward 0.022. Its intermediate collapse position (between GRASPrune at 18 and Standard at 60) reflects activation-norm-based selection that incidentally preserves some redundant structures with moderate norms.

\end{document}
